\section{參考實作}
\label{sec:design}

為驗證上述分析，本文實作 \system{}——
一套用於偵測故障節點並降低其路由權重的框架。
如第~\ref{sec:complementarity}~節所示，
在 P2C 下，偵測速度遠比具體的回應策略重要。
以下對 \system{} 的說明旨在描述評估中所採用的偵測機制。

\subsection{架構概覽}

\system{} 運行於控制平面，與 P2C 資料平面並行。
每個\emph{決策週期} $\Delta$（預設 0.5\,s），
執行三階段管線：\emph{監測} $\to$
\emph{偵測} $\to$ \emph{行動}。

在監測階段，系統從資料平面收集各節點指標：
請求離開間隔、佇列深度及全系統吞吐量。
將理論最大吞吐量與實測吞吐量之差定義為\emph{剩餘容量}，
並以指數移動平均（$\alpha = 0.3$）加以平滑。

在偵測階段，以離開間隔的 P10 百分位估計各節點服務時間。
根據 M/D/1 佇列理論，
此訊號不受負載影響（詳見第~\ref{sec:detection}~節）。
嚴重度定義為：
\begin{equation}
  \sigma_i = \max\left(0,\ 1 - \frac{\tbase}{d_i}\right)
  \label{eq:severity}
\end{equation}
其中 $d_i$ 為節點 $i$ 的估計服務時間。
當減速為 5 倍時 $\sigma = 0.8$；減速為 10 倍時 $\sigma = 0.9$。

在行動階段，根據嚴重度閾值選擇各節點策略，
並施加相應的流量調整。

\subsection{偵測：離開間隔估計}
\label{sec:detection}

偵測慢故障的核心難題在於區分真正的服務退化與暫態排隊效應。
在高負載（$\rho > 0.85$）下，
所有節點的延遲均會升高，使延遲式偵測器產生誤報。

本框架利用 M/D/1 佇列的一項性質：
當佇列持續處於忙碌狀態時（高利用率下多數時段皆如此），
離開間隔即為實際服務時間，不受佇列深度影響。
取離開間隔的 P10 百分位可濾除佇列清空時的空閒間隙，
從而逼近服務時間的下界。

此訊號具\emph{負載不變性}：
基底服務時間 $\tbase = 10$\,ms 的節點，
無論系統負載為 70\% 或 95\%，均產生 $d \approx 10$\,ms。
5 倍減速時，在任何負載水準下皆得 $d \approx 50$\,ms。

為避免暫態噪聲誤觸緩解機制，本框架要求\emph{持續確認}：
嚴重度須連續兩個決策週期超過最低閾值（$\sigma > 0.05$），
才會將該節點標記為故障。

\subsection{回應策略}
\label{sec:strategies}

如第~\ref{sec:complementarity}~節所示，
在 P2C 下具體的回應策略對 P99 的影響可忽略（$< 2$\,ms）。
本實作依嚴重度閾值
（$\theta_{\text{spec}} < \theta_{\text{shed}} < \theta_{\text{iso}}$）
選擇三種策略：
輕度退化時對沖（向另一健康節點發送副本），
中度退化時削減（依嚴重度降低權重 $w_i = \max(w_{\min},\, 1 - \sigma_i \cdot \text{spare\_factor})$），
嚴重退化時隔離（完全移除）。
隔離受容量邊界約束：
僅當隔離後剩餘有效利用率低於 $U$ 時方允許執行。

\subsection{負載感知閘控}

本框架的核心設計原則為：\emph{每項緩解行動皆有容量成本}。
對沖使請求負載加倍；
削減會將負載轉移至健康節點；
隔離則完全移除容量。
當系統負載偏高時，其成本可能超過效益，
反而形成緩解陷阱。

\system{} 以剩餘容量作為統一的閘控訊號，
同時管控三種策略：
\begin{itemize}
\item 對沖預算 $\propto$ 剩餘容量
  （剩餘容量為零時不進行對沖）。
\item 權重削減幅度 $\propto$ $\text{spare\_factor}$
  （剩餘容量 $< 0.08$ 時不削減）。
\item 當隔離後有效利用率 $\rhoeff > U$ 時，阻止隔離。
\end{itemize}

如此可確保 \system{} 優雅降級：
低負載時積極隔離以加速恢復；
高負載時保守削減以保留容量；
過載時則完全退避，避免適得其反。
